% Nouveau EGA V (cote 25): archive cote 25, site page 49 (PDF page 50).
% typescript. Draft compared: none.
% AI-transcribed by Claude Opus 5.5 (claude-opus-5-5, adaptive thinking, effort high), Message Batches API, batch ega25b (msgbatch_01A7n5YuXubUPNK1zsepi8RG), 2026-11-13; not yet checked by eye.
\scanpage{25/49}
\mspage{45}

toute composante irréductible \inserted{$X_i$} de dimension 1 \deleted{$X_i$} de X soit \xout{xxxbiex}
\inserted[typed]{et soit}
formée de points lisses de X \xout{xxxbien} réglée relativement à f, et que
pour tout point fermé \xout{de} x de X tel que $\dim_x X \geqslant 2$, on ait $\mathrm{prof}_x X \geqslant 2$,
(conditions automatiquement vérifiées si\unsure{.} X est \inserted{géom} normale et toutes ses
composantes irréductibles sont de dim$\geqslant$ 2).

b) Supposons X géom. normale , \xout{Pour} que l'ensemble \xout{Xxd} des $\hand{\zeta \in}\, \check{P}$
tels que $\phi_{\hand{\zeta}}$ soit $X_{k(\hand{\zeta})}$\hand{-}régulier et $Y_{\hand{\zeta}}$ \xout{soit} géom normale ait un complé-%
mentaire de codimension $\hand{\geqslant}$2, il f et s que \xout{pour} toute composante
irréductible $X_i$ de X de \xout{co}dimension$\hand{\leqslant}$2\unsure{\hand{,}} soit formée de point lisses de
X et soit réglée rel à f, et que de plus pour tout point fermé x de X
tel que $\dim_x X \hand{\geqslant} 3$, on ait $\mathrm{prof}_x X \hand{\geqslant}$ 3.

\xout{(J'aixoubliéxxdaxavantxa)xetxb)xxd'indiquerxl'hypothèse}
\xout{suivante,xtoujoursxvérifiéexenxcarx0xxpourxtoutexcompxirrxX$_i$xdexX$_i$,}
\xout{xxxxsupxsingx}
\xout{\unsure{g(X$_i$}xxxxxxxxxx)}

\uline{Remarque} 8.15. Dans 8.6. , 8.13. et 8.14. nous avons fait sur X
l'hypothèse $(S_k)$ (resp. $(R_k)$, resp\unsure{:} séparable, resp. \inserted[typed]{géom} normal ) que
nous désirions receuillir comme conclusion pour les sections hyperplanes
sauf \inserted{pour $\zeta$ dans} un ensemble exceptionnel de codimension $\geqslant$ 2. Cela ne restreint
pas la généralité (et il n'aurait pas coûté plus cher à vrai dire
de se dispenser de cette hypothèse préalable), car on voit aussitôt
à l'aide des résultats de par. 3.4. \inserted{et 5.12.} que si X ne satisfait pas à
l'hypothèse en question, donc (en vertu de par.5) s'il existe un point
\uline{fermé} x en lequel cette hypothèse est en défaut, alors pour tout $\hand{\zeta}$ tel
que $\phi_{\hand{\zeta}}$ soit $X_{k(\hand{\zeta})}$-régulier (condition n'éliminant qu'un ens de codim$\geqslant$2)
et tel que $x \hand{\in} Y_{\hand{\zeta}}$ (condition qui décrit un ensemble de codimension 1
exactement), $Y_{\hand{\zeta}}$ ne satisfait pas non plus à ladite hypothèse en x ,
\hand{donc} l'ensemble exceptionnel $Z \hand{\subset} \check{P}$ est de codimension 1 et non 2.
(J'ai peut-être charrié un peu dans le cas $(R_k)$, \over{o}{où} il faut encore
des conditions $(S_1)$ et d'équidimensionnalité peut-être.....).

\mspage{49}
